\subsection{Liaison iono-covalente}
\label{subsection:iono-covalent-bond}


\setcounter{equation}{0}
\setcounter{Q}{0}

Le dioxyde de cobalt \ce{Li_xCoO2} est étudié à l'aide d'une courbe de charge/décharge à une vitesse $C/24$ sur une plage de potentiel de $3.6$ à \SI{4.85}{\volt/{Li^+ / Li}}. La courbe de charge/décharge suivante est obtenue :

\begin{center}
\begin{minipage}[]{0.6\textwidth}
\centering
\includegraphics[height=5cm]{Ex18/Fig01}
\end{minipage}\begin{minipage}[]{0.3\textwidth}
\captionof{figure}{Courbes de charge/décharge de \ce{Li_xCoO2} à un courant électrique valant $C/24$.}
\label{fig:18-01}
\end{minipage}
\end{center}

La capacité spécifique expérimentale de la cathode vaut : $C = 140 ~ \si{\milli\ampere\hour\per\gram}$. Il est courant d'exprimer le courant en l'indiquant sous la forme :
\begin{align*}
\text{courant (massique)} = \frac{\text{capacité (spécifique massique) de la cellule}}{\text{durée de la décharge}}
\end{align*}

\Q Quel est le sens du terme $C/24$ ? Trouver le courant appliqué par gramme de matériau.

\solution{
Le courant $C/24$ signifie que l'on a déchargé le matériau à vitesse constante de telle sorte que la batterie atteigne l'état $SoC=0 ~ \si{\percent}$ en \SI{24}{\hour}. Par conséquent, la vitesse par gramme de matériau, notée $I_m$, vaut : 
\begin{align*}
I_m = \frac{C}{\Delta t} \stackrel{AN}{=} \frac{140}{24} \approx 5.8 ~ \si{\milli\ampere\per\gram}. 
\end{align*}
\faWarning ~ Cette notation n'est pas abordée dans le cours. Elle est précisée dans cet exercice, question de dire qu'elle ne soit pas étrangère aux étudiants si elle est rencontrée dans une publication. On trouve des tonnes de discussion en ligne sur cette notation, ils n'auront aucun soucis à se débrouiller avec si elle est rencontrée. En partiel, la définition de la notation sera donnée dans les données si nécessaire.
}

\Q Donner l'évolution des degrés d'oxydation de l'oxygène et du cobalt lors de la charge dans un modèle de liaison \SI{100}{\percent} ionique.
\solution{En utilisant l'équation de conservation de la charge dans un composé $\ce{Li_x Co O2}$: 
\begin{align*}
x \cdot no (\ce{Li}) + no (\ce{Co}) + 2 \cdot no (\ce{O}) = 0
\end{align*}
On en déduit les résultats suivants pour les motifs limites : 
\begin{center}
\begin{tabular}{ccc}
\toprule
SoC & \SI{0}{\percent} & \SI{100}{\percent} \\
Fraction lithium x & 1 & 0 \\
\midrule
$no (\ce{Co})$ & $+III$ & $+IV$ \\
$no (\ce{O})$ & $-II$ & $-II$ \\
\bottomrule
\end{tabular}
\end{center}
}

On rappelle que la structure de \ce{LiCoO2} est lamellaire et que cette structure est décrite par une maille hexagonale. La structure lithiée a été suivie par DRX. On a constaté que pour $x > 0.5$, la structure ne subit aucun changement dans ses paramètres de maille. Cependant, pour $x<0.5$, on constate que le rapport $c/a$ de la maille hexagonale diminue.

\Q En vous appuyant sur la structure lamellaire, que peut traduire microscopiquement une diminution du rapport $c/a$ ?
\solution{Une diminution du rapport $c/a$ traduit le fait que l'espace entre les plaques diminue. Or, on sait que l'origine de la distance entre les plaques est électrostatique : c'est la répulsion entre les charges partielles négatives des oxygènes qui se font face.}

On a vu que le cobalt est entouré d'oxygène dans un environnement octaédrique, de coordinence [6]. 

\Q En vous appuyant sur le modèle du champ cristallin, rappeler, en le justifiant, le diagramme d'orbitale du bloc $d$ du cobalt $+III$.

Il a été expérimentalement observé que lors du passage de $x = 1$ à $x = 0$, une transition de spin, d'une structure bas spin à une structure haut spin. La mesure expérimentale de l'évolution des longueurs de liaison a montré que selon la direction $z$, il y a un allongement significatif des octaèdres de cobalt conjointe à une compression dans la direction $xy$.

\Q Cette observation est-elle en accord avec les mesures de DRX précédente ?

\solution{La structure DRX nous donne un allongement selon la direction $c$ et une compression selon la direction $a$. Cette observation est en accord avec les mesures réalisées sur les longueurs de liaisons.}


\Q \`{A} partir du diagramme d'énergie de l'octaèdre prédit par le champ cristallin, proposer une structure électronique qui rende compte du nouvel état de spin.

\emph{NB : On va supposer que la déformation est suffisante pour que le niveau $d z^2$ devienne plus stable que le niveau $d_{x^2 - y^2}$}

\solution{
En cas d'expansion de l'octaèdre dans la direction $z$ et une contraction dans la direction $xy$, on s'attend à une :
\begin{itemize}
\item diminution de la répulsion entre les électrons du métal et des ligands dans la direction $z$, donc une forte stabilisation en énergie de la $d_{z^2}$ et dans une moindre mesure des $d_{xz} $ et $d_{yz}$ ;
\item augmentation de la répulsion dans le plan $xy$ entre les électrons du métal et des ligands, donc une forte déstabilisation en énergie de la $d_{x^2 - y^2}$ et dans une moindre mesure de la $d_{xy}$.
\end{itemize}
\begin{center}
\begin{minipage}{0.5\textwidth}
\includegraphics[width=\textwidth]{Ex39/Fig03a}
\end{minipage}\begin{minipage}{0.46\textwidth}
orbitales {\color{red} -} destabilisées \\
orbitales {\color{blue} -} stabilisées
\end{minipage}
\end{center}
L'octaèdre déformé correspond à la structure délithié $\ce{CoO2}$ et l'octaèdre parfait à la structure lithié \ce{LiCoO2}. Par conséquent, on s'attend à avoir une structure $d^6$ champ fort, bas spin pour l'octaèdre parfait, et une structure $d^5$ pour l'octaèdre déformé.
\begin{center}
\includegraphics[width=0.5\textwidth]{Ex39/Fig03b}
\end{center}
}


\Q En vous appuyant sur une interaction selon l'axe $z$ entre une orbitale $d_{z^2}$ et une orbitale $2p_z$, montrer la différence entre l'oxydation dans un modèle purement ionique et un modèle iono-covalent.
\solution{
Dans un modèle ionique, l'écart en énergie des orbitales est suffisamment grand pour que les orbitales du métal soient les seules à perdre des électrons : 
\begin{center}
\begin{minipage}{0.48\textwidth}
\centering
\includegraphics[width=\textwidth]{Ex39/Fig01a}
Avant oxydation
\end{minipage}\hfill \begin{minipage}{0.48\textwidth}
\centering
\includegraphics[width=\textwidth]{Ex39/Fig01b}
Après oxydation
\end{minipage}
\end{center}
Dans un modèle iono-covalent, l'orbitale oxydée présente de la densité électronique sur le métal, mais aussi, dans une moindre mesure, sur l'oxygène :
\begin{center} 
\begin{minipage}{0.48\textwidth}
\centering
\includegraphics[width=\textwidth]{Ex39/Fig02a}
Avant oxydation
\end{minipage}\hfill \begin{minipage}{0.48\textwidth}
\centering
\includegraphics[width=\textwidth]{Ex39/Fig02b}
Après oxydation
\end{minipage}\end{center}}

\Q Que peut-on dire de l'évolution de la charge partielle aux oxygène dans un modèle ionocovalent ?

\solution{
Le diagramme d'interaction ionocovalent montre qu'après oxydation, les électrons occupant l'orbitale moléculaire antiliante  $d_{z^2} - p_z$ sont perdus. Puisqu'il y a de la densité électronique sur l'oxygène, c'est-à-dire que l'électron avait une probabilité non nul d'être sur l'oxygène, l'oxydation conduit à une diminution de la charge partielle portée par l'oxygène.
}

